\documentclass[11pt]{article}

\usepackage[utf8]{inputenc}
\usepackage[T1]{fontenc}
\usepackage{amsmath, amssymb}
\usepackage{geometry}
\usepackage{hyperref}
\usepackage{enumitem}
\geometry{letterpaper, margin=1in}

\title{\textbf{Working Draft: Paper 1 \\ The Bipolar Graviton Condensate as an Effective Non-Singular Core Model}}
\author{Stephen Breighner \\ \small Independent Researcher}
\date{March 2026}

\begin{document}
\maketitle

\begin{abstract}
This paper develops the conservative branch of the Bipolar Graviton Condensate program. The aim is narrow: treat the condensate as an effective description of a finite-radius compact core, keep the long-range gravitational baseline close to the standard massless picture, and defer any explicit matter-to-condensate conversion mechanism to later work. Within that limited scope, an effective potential with a short-distance support term, a mid-range binding term, and a long-range gravitational term admits a stable radius and a modified fringe velocity profile. The result is not presented as a derivation from fundamental quantum gravity, but as a consistency model for how a non-singular core and a small deviation from a pure Schwarzschild profile could coexist.
\end{abstract}

\section{Introduction}
General Relativity describes gravitational collapse with extraordinary success at large scales, but the classical Schwarzschild picture still drives the interior solution toward a singular point. The present paper does not attempt to resolve that problem from first principles. Instead, it adopts a more limited objective: ask whether a compact gravitational end-state can be modeled as a finite-radius effective core without immediately contradicting the long-range behavior of gravity.

The Bipolar Graviton Condensate is used here as a phenomenological ansatz rather than a completed microscopic theory. In this conservative branch, the key claims are deliberately restricted:
\begin{itemize}[leftmargin=1.25em]
  \item the condensate is treated as an effective phase or description, not a proven fundamental particle ontology;
  \item the long-range gravitational sector is kept close to the standard massless baseline;
  \item the short-distance support term is introduced explicitly as a regularization term, not yet as a derived interaction;
  \item observable interest is focused on two outputs of the model: a finite stable radius and a modified fringe velocity profile.
\end{itemize}

This framing separates the first paper from more speculative branches. Questions about literal matter-to-graviton transfer, transfer-field dynamics, or extended constituent structure are deferred to later papers.

\section{Conservative Assumptions}
The paper is organized around an effective one-dimensional radial model. The separation coordinate $r$ represents a characteristic radial scale of the hypothesized compact core. It is not claimed here that $r$ is the radius of a sharply bounded object, only that it measures the scale at which the modeled interactions balance.

The conservative assumptions are:
\begin{enumerate}[leftmargin=1.4em]
  \item gravity retains its standard long-range form to leading order;
  \item the condensate can be modeled by an effective potential without committing to a full microscopic derivation;
  \item short-distance support may be represented phenomenologically by a steep repulsive or pressure-like term;
  \item the theory is judged in this paper by internal consistency and by the form of its radial profiles, not by a claim of completed quantum-gravity derivation.
\end{enumerate}

This is therefore a controlled narrowing of scope. The model is intended to answer the question: if an effective bipolar condensate phase exists, what minimum radial structure follows from that assumption?

\section{Effective Potential Ansatz}
In normalized units with $G=1$ and $M=1$, the effective potential is written as
\[
V(r) = \frac{A}{r^{12}} - \frac{B}{r^3} - \frac{1}{r}.
\]
The three terms are interpreted as follows:
\begin{itemize}[leftmargin=1.25em]
  \item $A/r^{12}$ is a short-distance support term that prevents collapse to $r=0$ within the model;
  \item $-B/r^3$ is a mid-range bipolar binding term;
  \item $-1/r$ preserves the usual long-range gravitational attraction in normalized form.
\end{itemize}

The associated radial force is
\[
F(r) = -\frac{dV}{dr} = \frac{12A}{r^{13}} - \frac{3B}{r^4} - \frac{1}{r^2}.
\]
A stable radius $r_*$ is defined by two conditions:
\[
F(r_*) = 0, \qquad \frac{d^2V}{dr^2}\Big|_{r_*} > 0.
\]
The second derivative test distinguishes a stable minimum from an unstable turning point.

A useful limiting case is obtained by setting $A=0$:
\[
V(r) = -\frac{B}{r^3} - \frac{1}{r},
\]
which yields
\[
F(r) = -\frac{3B}{r^4} - \frac{1}{r^2} < 0 \qquad (r>0).
\]
In other words, once the support term is removed, the model has no finite static equilibrium radius. This observation matters because it makes explicit where singularity avoidance enters the present construction: it does not arise automatically from the attractive terms alone.

\section{Interpretation of the Support Term}
The support term is the most important caution point in the paper. It is included because the model requires some short-distance mechanism that resists unlimited collapse. At this stage, the term should be read as an effective placeholder for short-distance support rather than as a final microscopic law.

Several broad interpretations remain open for later work:
\begin{itemize}[leftmargin=1.25em]
  \item quantum or gradient pressure associated with a condensate field;
  \item repulsive self-interaction in an underlying bipolar sector;
  \item finite-size or non-overlap effects of an extended constituent picture;
  \item a modified stress contribution that regularizes the core at high density.
\end{itemize}

This paper does not choose among those possibilities. The only claim made here is conditional: if a short-distance support mechanism of this general kind exists, then the effective model admits a finite-radius stable core.

\section{Fringe Velocity Curve}
To compare the modeled condensate against a Schwarzschild-like baseline, define the orbital-velocity proxy
\[
v(r) = \sqrt{r\,|F(r)|}.
\]
For the normalized Schwarzschild baseline,
\[
F_{\mathrm{Sch}}(r) = -\frac{1}{r^2},
\qquad
v_{\mathrm{Sch}}(r) = \frac{1}{\sqrt{r}}.
\]
For the bipolar condensate ansatz,
\[
v_{\mathrm{BGC}}(r) = \sqrt{r\left|\frac{12A}{r^{13}} - \frac{3B}{r^4} - \frac{1}{r^2}\right|}.
\]
The short-distance support term dominates only very near the core, while the mid-range $1/r^3$ binding term can alter the profile at larger radii than the support wall itself. In the intended interpretation, this does not replace the asymptotic gravitational falloff. Instead, it produces a mild fringe-region deformation of the velocity curve relative to the pure Schwarzschild case.

That distinction is central. The conservative branch is not arguing for a wholesale replacement of long-range gravity. It argues only that a finite-radius core with an additional mid-range structure can leave a residual signature in the outer profile.

\section{What This Paper Claims and Does Not Claim}
The present draft supports the following claims:
\begin{itemize}[leftmargin=1.25em]
  \item an effective radial potential can be written down that remains close to the standard long-range gravitational form while avoiding $r=0$ collapse inside the model;
  \item the same effective model produces a finite stable radius when the support term is present;
  \item the model predicts a fringe velocity profile that differs slightly from the pure Schwarzschild baseline.
\end{itemize}

The draft does \textit{not} claim the following:
\begin{itemize}[leftmargin=1.25em]
  \item that the microscopic origin of the support term is already known;
  \item that ordinary matter has been shown to literally convert into gravitons;
  \item that the model is a full solution of Einstein's equations with a derived stress-energy tensor;
  \item that the existence of a stable minimum in the effective potential is experimental proof of a real bipolar graviton condensate.
\end{itemize}

This boundary is important for the paper's credibility. The first paper remains strongest when it is explicit about being a controlled effective model.

\section{Relation to Later Papers}
By construction, this paper leaves two larger branches for later treatment.

First, a transfer-field paper can ask how collapsing matter enters the condensate description. That branch would need trigger conditions, continuity equations, and conservation-law accounting. Second, an extended-constituent paper can ask whether the short-distance support term arises from a massive or structured bipolar constituent picture. That branch would need to revisit both the internal support mechanism and the large-scale force law.

The present paper is therefore intentionally incomplete in a productive way. Its job is to establish the narrowest defensible first step before moving to more speculative mechanisms.

\section{Conclusion}
The Bipolar Graviton Condensate can be introduced conservatively as an effective non-singular core model rather than as a complete microscopic theory. In that limited role, the model provides two useful outputs: a finite stable radius and a small deviation in the fringe velocity profile relative to a pure Schwarzschild baseline. The price of this result is clear and should remain explicit: a short-distance support term is required, and its origin is not yet derived in this paper.

That tradeoff is acceptable for a first paper provided the scope remains disciplined. The effective model does not prove that nature contains bipolar gravitons, nor does it solve quantum gravity. It does, however, show that a compact, non-singular radial structure with mild outer-profile consequences can be formulated in a mathematically transparent way. That is enough to justify a first paper focused on model consistency, while leaving conversion mechanisms and constituent structure to later branches.

\vspace{1em}
\noindent\textbf{References}\\
\small
Einstein, A. (1915). \textit{The Field Equations of Gravitation}.\\
Saldanha, P. L., et al. (2026). \textit{Repulsive Gravitational Force as a Witness of the Quantum Nature of Gravity}.\\
Breighner, S. (2026). \textit{The Bipolar Graviton Condensate}. Current project manuscript.

\end{document}
